\documentclass[11pt]{article}
\usepackage[margin=1in]{geometry}
\usepackage[hidelinks]{hyperref}
\usepackage{tabularx}

\begin{document}

\begin{titlepage}
\centering
{\Large George Washington University Data Science Program\par}
\vspace{1cm}
{\large Capstone Report --- Fall 2026\par}
\vspace{1.5cm}
{\huge\bfseries Exploring Applications for the Montandon Global Crisis Data Bank\par}
\vspace{2cm}
{\large Jeongmin An, Jehan Bugli, and Aidan Carlisle\par}
\vspace{1cm}
{\large Advisor: Amir Jafari\par}
\vfill
{\large Draft --- September 2026\par}
\end{titlepage}

\tableofcontents
\newpage

\section{Introduction}

This report documents a project exploring applications of the Montandon Global Crisis
Data Bank. The project combines disaster and response records with external news,
geospatial context, text embeddings, and network analysis.

\section{Problem Statement}

Montandon contains rich crisis and response information, but connecting related records
and exploring their relationships can be difficult. This project investigates whether a
graph-based representation and a simple query interface can make that information easier
to explore. Results and evaluation criteria will be added as the project is scoped.

\section{Related Work}

Crisis risk reduction depends on understanding disaster risk, preparedness, and recovery
\cite{undrr2015sendai}. Research on emergency social-media messages
includes filtering, classification, and summarization \cite{imran2015processing}; this
project instead starts with the curated Montandon data bank and uses external articles as
additional context. STAC provides a common structure for spatiotemporal metadata
\cite{ogc2025stac}, which is relevant because Montandon uses a modified STAC approach.

Adding weather and news to Montandon requires matching records across sources.
GDIS does this for EM-DAT, adding administrative boundaries and coordinates while
keeping the original event identifiers \cite{rosvold2021gdis}. Missing locations
and later EM-DAT updates can leave gaps or mismatches in the joined data.
EM-DAT itself combines disaster loss reports from multiple sources and documents
how it selects and checks the figures \cite{delforge2025emdat}. Its uneven
reporting across countries and years also limits comparisons.

News can provide both impact figures and relationships between responders.
Comfort and Haase reconstructed Hurricane Katrina's response network from news
reports, then examined organizational interactions after communications were
damaged \cite{comfort2006communication}. Our graph has some similarities
in spirit, using recorded relationships, shared attributes, and more
to connect disaster and IFRC response records.
Li et al. used LLMs and conventional NLP to extract climate impact
records from online reports \cite{li2024climateimpacts}. They compared three
models against an existing database and found that accuracy varied by field.
Our news pipeline currently retrieves and ranks articles; extracting potential
network relationships extracting facts from them would extend that work.

The graph also gives an LLM access to records beyond its training data.
GraphRAG builds its graph from documents and summarizes connected groups of
entities, reporting more comprehensive answers to wide-ranging questions than
a conventional RAG baseline \cite{edge2024graphrag}. We build from structured
records and retrieve through predefined Cypher queries. Generating queries directly
is another option; Ozsoy et al. assembled 44,387 question-to-Cypher examples and
improved reference query agreement through fine-tuning \cite{ozsoy2025text2cypher}.
Our assistant chooses a tool and its arguments, leaving the Cypher fixed.

That interface can be evaluated using the tasks like those in the Berkeley Function
Calling Leaderboard: choosing functions, supplying arguments, making several
calls, and recognizing when no call is needed \cite{patil2025bfcl}. The benchmark
finds that models still struggle with longer tasks despite strong performance on
single calls. We plan to test some of these behaviours on Montandon and IFRC GO questions
and check the resulting answers against the records returned.

\section{Solution and Methodology}

This section describes the data sources and will be expanded with feature generation,
graph construction, and the query setup.


\subsection{Data Sources}

This project combines disaster records obtained through Montandon with
selected external data sources. Montandon records were obtained through its
API using a token issued through IFRC GO. Prepared datasets used to build
the graph are stored in S3 and loaded into Neo4j. The table summarizes
each source's coverage,  date range in our project files,
licence or use terms, and graph mapping.

\begin{table}[htbp]
    \centering
    \footnotesize
    \begin{tabularx}{\textwidth}{@{}p{2.3cm}X p{2.4cm}p{2.5cm}X@{}}
    \hline
    Source (origin) & What it covers & Date range in project files &
    Access / use terms &  How it enters the graph \\
    \hline

    EM-DAT (Montandon) &
    Disaster events and reported impacts &
    1900-01-01--2026-08-28 &
    Montandon API token (issued through IFRC GO) &
    \texttt{Event} and \texttt{Impact} nodes store the records;
    \texttt{IMPACT\_OF} links impacts to events \\
    \hline

    GDACS (Montandon) &
    Disaster alerts and event information &
    2000-01-01--2026-09-11 &
    Montandon API token (issued through IFRC GO); Free to use with acknowledgement of GDACS &
    \texttt{Event} nodes store the records and link to country,
    hazard, and date nodes \\
    \hline

    IFRC event collection (Montandon) &
    Disaster-response events and impacts &
    Events: 1970-01-01--2026-08-26;
    impacts: 1997-11-10--2026-08-10 &
    Montandon API token (issued through IFRC GO) &
    \texttt{Event} and \texttt{Impact} nodes store the records;
    \texttt{IMPACT\_OF} links impacts to events \\
    \hline

    IFRC GO (external) &
    Emergency events and appeals &
    Events: 1919-01-01--2026-09-11;
    appeals: 1919-01-01--2026-09-02 &
    Public IFRC GO API &
    \texttt{GOEvent} and \texttt{GOAppeal} nodes;
    \texttt{IN\_COUNTRY} and \texttt{FOR\_EVENT} relationships \\
    \hline

    NASA POWER (external) &
    Daily temperature, precipitation, and wind data for locations
    associated with Montandon records &
    2021-09-15--2026-09-13 &
    Free access; cite NASA POWER &
    Daily records are matched by \texttt{item\_id} and summarized
    as weather properties on \texttt{MontandonItem} nodes \\
    \hline

    Media Cloud (external) &
    News article metadata, with article text retrieved separately
    from publisher websites &
    \textemdash &
    Media Cloud account and API key &
    \texttt{NewsArticle} nodes store article information and link
    to \texttt{Event} nodes through \texttt{VERIFIED\_FOR} \\
    \hline

    Geospatial data (external) &
    &
    &
    &
    \\
    \hline
    \end{tabularx}
    \end{table}

    Source-access and citation information is documented by
    \href{https://montandondata.org/get-started.html}{Montandon},
    \href{https://doc.emdat.be/docs/data-accessibility/}{EM-DAT},
    \href{https://www.gdacs.org/Documents/2025/GDACS_API_quickstart_v2.pdf}{GDACS},
    \href{https://go.ifrc.org/terms-and-conditions}{IFRC GO},
    \href{https://power.larc.nasa.gov/docs/referencing/}{NASA POWER},
    \href{https://www.mediacloud.org/legal/media-cloud-terms-of-use}{Media Cloud}.




\section{Results and Discussion}

This section will report the project’s results, figures, and limitations.

\section{Conclusion}

This section will summarize findings and recommended next steps.

\bibliographystyle{IEEEtran}
\bibliography{references}

\end{document}
